%% Introduction — \neven{} Paper
%% Revised: reduced "and" repetition, added cost barrier, R4XCL origin, broader scope

\section{Introduction}
\label{sec:introduction}

\subsection{The Spreadsheet Dominance Problem}

Microsoft Excel is estimated to serve over 750 million users worldwide \citep{Microsoft2023}, making it the most widely deployed data analysis tool across scientific research, corporate finance, public policy, and education. Despite this dominance, studies have documented that approximately 88\% of spreadsheets contain errors \citep{Panko2016}. High-profile failures underscore the risks: the London Whale trading loss of \$6.2 billion was attributed in part to spreadsheet errors \citep{JPMorgan2013}; the truncation of COVID-19 case data in England resulted from Excel's row limits \citep{Kelion2020}; influential economic policy conclusions were shown to rest partly on spreadsheet coding mistakes \citep{Herndon2014}.

In scientific research, the reproducibility crisis has been extensively documented \citep{Baker2016, Ioannidis2005}. Spreadsheet-based analyses are particularly vulnerable because they lack version control, audit trails, and formal testing mechanisms \citep{Peng2011}. Yet practitioners continue to use Excel because of its accessibility, immediate visual feedback, and low barrier to entry \citep{Scaffidi2005}. The tool is not the problem---the absence of computational infrastructure beneath it is.

\subsection{The Programming Expertise Barrier}

Scientific computing environments such as R \citep{RCoreTeam2024}, Julia \citep{Bezanson2017}, and Python \citep{VanRossum2009} provide rigorous, reproducible analytical capabilities spanning statistics, numerical mathematics, machine learning, optimization, and visualization. However, these tools require programming expertise that represents a significant barrier for domain experts in economics, biology, social sciences, and business analytics. \citet{Muenchen2012} documented that the learning curve of R remains a primary obstacle to adoption among applied researchers. Julia, while offering superior numerical performance \citep{Bezanson2017}, presents an even steeper adoption curve due to its smaller ecosystem.

This creates a fundamental tension: the tools that are accessible lack computational rigor, while the tools that provide rigor lack accessibility. Previous attempts to bridge this gap---including RExcel \citep{Baier2011}, BERT \citep{BERT2018}, xlwings \citep{xlwings2023}, PyXLL \citep{PyXLL2023}, and Microsoft's Python in Excel \citep{MicrosoftPython2023}---have addressed only partial aspects of the problem. Each typically supports a single language, lacks security isolation or interactive visualization, and several impose commercial licensing costs that exclude academic institutions with limited budgets. None offers a unified, open-source platform integrating multiple computational ecosystems simultaneously.

\subsection{Contributions}

This paper presents \neven{}, an open-source polyglot infrastructure that addresses the computational accessibility gap through the following contributions:

\begin{enumerate}
    \item \textbf{Simultaneous polyglot integration:} A C++17 architecture that unifies R~4.4.1, Julia~1.12.6, Python~3.13, Quarto, and JavaScript visualization libraries into a single platform accessible from Excel formulas. Each of these ecosystems is individually capable of transforming an entire discipline; \neven{} places them all in the hands of non-programmer users through a formula-level interface.

    \item \textbf{Unified multi-format viewer:} A WebView2 (Chromium-based) engine that renders interactive Plotly charts, D3.js visualizations, rpivotTable pivot tables, Leaflet geographic maps, Markdown with LaTeX mathematics, PDF documents, Word files (via Pandoc), and Pluto.jl reactive notebooks---all through a single Excel formula (\texttt{=NEVEN.v()}).

    \item \textbf{Process-isolated fault-tolerant architecture:} Language runtimes execute in separate OS processes communicating via Named Pipes with Protocol Buffers serialization. A crash in any runtime does not propagate to the host application. A security sandbox blocks over 30 dangerous patterns per language, validated by 109 automated tests.

    \item \textbf{User extensibility:} A file-based extension model where users add custom R, Julia, or Python functions by placing script files in a directory. Functions are automatically discovered, registered with Excel, and exposed as native formulas---enabling domain-specific teams to build specialized libraries without modifying the core system.

    \item \textbf{Reproducible publishing:} Integration with Quarto for generating publication-quality reports (HTML, PDF, Word) that embed R/Julia/Python computations alongside narrative text, citations, and cross-references---completing the pipeline from raw data to peer-reviewed publication within the Excel environment.

    \item \textbf{Empirical validation:} Verification of computational outputs against the established econometric benchmarks of \citet{Wooldridge2016}, demonstrating numerical correctness across regression, panel data, and time series implementations.

    \item \textbf{Open-source community model:} Released under the GNU General Public License v3, \neven{} is explicitly designed so that every user can extend its capabilities by contributing functions in any supported language, growing the platform through shared knowledge. The project evolved from R4XCL \citep{BonillaR4XCL2023}, the author's master's thesis work at the Universidad de Costa Rica, which originally connected R with Excel as a single-language bridge.
\end{enumerate}

\subsection{Paper Organization}

The remainder of this paper is organized as follows. Section~\ref{sec:architecture} describes the system architecture, including the IPC protocol, security model, and visualization subsystem. Section~\ref{sec:viewer_extensibility} presents the unified viewer, user extensibility model, and reproducible publishing pipeline. Section~\ref{sec:case_studies} validates the system through econometric benchmarks and an enterprise analytics scenario. Section~\ref{sec:discussion} discusses performance characteristics, comparison with alternatives, and limitations. Section~\ref{sec:conclusion} concludes with future work directions.
